\chapter{Visuomotor Policies}

\section*{학습 목표}

이 장은 언어가 로봇 policy의 중심에 들어오기 전부터 존재했던 image-to-action 계보를 다룬다. VLA는 language-conditioned policy이지만, 실제 action을 생성하고 실행하는 하위 문제는 여전히 visuomotor control의 문제이다. 카메라가 본 장면, proprioception, gripper state, 과거 관측, task condition을 이용해 다음 action을 내는 구조가 오늘날 VLA의 직접적인 전신이다.

\begin{enumerate}
  \item visuomotor policy를 observation-to-action mapping으로 정의한다.
  \item visual observation이 partial, delayed, viewpoint-dependent state라는 점을 설명한다.
  \item temporal memory, recurrent policy, attention, transformer policy가 왜 필요한지 정리한다.
  \item real-robot grasping, pushing, reaching, insertion에서 필요한 evaluation metadata를 정의한다.
\end{enumerate}

\section{From vision to action}

Visuomotor policy는 시각 관측과 로봇 상태를 받아 action을 출력하는 policy이다 \cite{levine2016gps,ebert2018visualforesight,florence2021transporter,cliport2022}.
\begin{equation}
  \act_t
  =
  \pi_{\theta}(\bm{I}_t,\conf_t,\bm{g}_t),
  \label{eq:visuomotor-policy-basic}
\end{equation}
여기서 $\bm{I}_t$는 camera image, $\conf_t$는 robot configuration 또는 proprioception, $\bm{g}_t$는 goal representation이다. goal은 목표 이미지, one-hot task label, language embedding, keypoint, affordance map 등 여러 형태가 될 수 있다.

VLA 관점에서 보면 language가 추가되기 전의 policy는 다음과 같은 구조였다.
\begin{equation}
  \act_t
  =
  \pi_{\theta}(\bm{I}_t,\conf_t,\bm{z}^{task}),
  \label{eq:task-conditioned-visuomotor}
\end{equation}
여기서 $\bm{z}^{task}$는 task id, goal image, desired pose, 또는 learned task embedding이다. VLA는 이 자리에 language instruction $\task$ 또는 VLM embedding을 넣는다.

\begin{definitionbox}{Visuomotor policy}
Visuomotor policy는 visual observation과 robot state를 조건으로 하여 robot action을 생성하는 policy이다. VLA는 여기에 language-conditioned task semantics를 추가한 형태로 볼 수 있지만, visual feedback, action timing, controller interface, distribution shift 문제는 그대로 남는다.
\end{definitionbox}

\section{Visual observation is not state}

카메라 영상은 풍부하지만 완전한 state가 아니다. 물체가 보인다고 해서 정확한 6-DoF pose, mass distribution, friction, contact state, occluded obstacle을 아는 것은 아니다. Chapter 2에서 다룬 observation/state 구분은 visuomotor policy에서 가장 직접적으로 나타난다.

visual encoder를 $\phi_{\theta}$라 하면,
\begin{equation}
  \bm{h}_t
  =
  \phi_{\theta}(\bm{I}_t),
  \qquad
  \act_t
  =
  \pi_{\theta}(\bm{h}_t,\conf_t).
  \label{eq:visual-encoder-policy}
\end{equation}
하지만 $\bm{h}_t$가 충분한 state representation인지는 task와 sensor setup에 따라 다르다. camera calibration, viewpoint, lighting, occlusion, motion blur, rolling shutter, depth noise가 모두 action 품질에 영향을 준다.

\begin{warningbox}{Image-to-action의 숨은 가정}
Image-to-action policy는 image가 task-relevant state를 충분히 포함한다는 가정을 둔다. 이 가정이 깨지면 policy는 시각적으로 그럴듯한 action을 내더라도 contact, force, hidden object, timing 문제에서 실패한다.
\end{warningbox}

\section{Closed-loop visual feedback}

open-loop trajectory를 한 번 생성해서 실행하는 방식은 작은 perturbation에 약하다. Visuomotor policy의 강점은 매 step 관측을 다시 보고 action을 수정할 수 있다는 점이다.
\begin{equation}
  \act_t = \pi_{\theta}(\obs_t,\conf_t),
  \qquad
  \state_{t+1}=f(\state_t,\controller(\act_t,\hat{\state}_t),\bm{w}_t).
  \label{eq:closed-loop-visuomotor}
\end{equation}
이 closed-loop 구조는 reaching, grasping, pushing처럼 visual correction이 필요한 task에서 중요하다.

Visual servoing의 고전적 형태는 feature error를 줄이는 control이다. image feature를 $\bm{y}$, target feature를 $\bm{y}^{ref}$라 하면,
\begin{equation}
  \dot{\conf}
  =
  -K J_{img}(\conf)^{\dagger}(\bm{y}-\bm{y}^{ref}),
  \label{eq:visual-servoing}
\end{equation}
여기서 $J_{img}$는 image feature velocity와 joint velocity를 연결하는 interaction matrix이다. Learned visuomotor policy는 이 mapping을 명시적으로 쓰지 않을 수 있지만, 실제로는 visual error를 action correction으로 바꾸는 문제를 학습한다.

\section{Temporal memory}

단일 image는 속도, 접촉 전후 변화, occluded state를 충분히 알려주지 않는다. 따라서 policy는 history를 사용할 수 있다.
\begin{equation}
  \act_t
  =
  \pi_{\theta}(\obs_{t-H:t},\conf_{t-H:t},\bm{g}_t).
  \label{eq:history-policy}
\end{equation}
recurrent policy는 hidden state를 갱신한다.
\begin{equation}
  \bm{m}_t
  =
  F_{\theta}(\bm{m}_{t-1},\obs_t,\conf_t),
  \qquad
  \act_t
  =
  G_{\theta}(\bm{m}_t,\bm{g}_t).
  \label{eq:recurrent-visuomotor}
\end{equation}

VLA에서도 memory는 optional feature가 아니다. language instruction이 길거나 task가 multi-stage일 때, policy는 현재 frame뿐 아니라 이전 subgoal, grasp 시도, contact event, failure signal을 기억해야 한다.

\section{Attention and transformer policies}

CNN 기반 encoder는 image feature를 만들고, MLP 또는 recurrent head가 action을 출력하는 구조가 오래 쓰였다. Transformer policy는 image patch, proprioception token, action history, task token을 하나의 sequence로 처리할 수 있다.
\begin{equation}
  \bm{Z}_t
  =
  [\bm{z}^{img}_{1:P},\bm{z}^{prop}_t,\bm{z}^{task},\bm{z}^{hist}_{t-H:t}],
  \qquad
  \bm{H}_t = \mathrm{Transformer}_{\theta}(\bm{Z}_t).
  \label{eq:transformer-policy-tokens}
\end{equation}
action head는 transformer output에서 action distribution을 만든다.
\begin{equation}
  p_{\theta}(\act_t\mid \obs_{\le t},\task)
  =
  \mathrm{ActionHead}_{\theta}(\bm{H}_t).
  \label{eq:transformer-action-head}
\end{equation}

이 구조는 VLA와 자연스럽게 연결된다. language token과 image token을 같은 transformer에 넣을 수 있고, robot action을 token, continuous vector, chunk, diffusion target으로 출력할 수 있다. 그러나 transformer backbone이 들어왔다고 해서 controller rate, latency, contact feedback 문제가 사라지는 것은 아니다.

\section{Real-robot grasping}

Grasping은 visuomotor policy의 대표적 testbed이다. image에서 grasp point, approach direction, gripper width, lift action을 예측하는 구조는 다음처럼 쓸 수 있다.
\begin{equation}
  \act_t
  =
  (\bm{p}^{grasp}_t,\bm{R}^{grasp}_t,w_t,\alpha_t),
  \label{eq:grasp-action}
\end{equation}
여기서 $\bm{p}^{grasp}$는 grasp position, $\bm{R}^{grasp}$는 grasp orientation, $w$는 gripper width, $\alpha$는 execute/open/close 같은 discrete mode이다.

grasp policy의 실패는 여러 계층에서 발생한다.

\begin{itemize}
  \item perception failure: object segmentation 또는 pose 추정 실패
  \item action failure: grasp pose가 물체 geometry와 맞지 않음
  \item control failure: approach trajectory나 gripper timing 문제
  \item contact failure: slip, insufficient normal force, collision
  \item recovery failure: 실패 후 재파지하지 못함
\end{itemize}

따라서 grasp success rate만으로는 policy를 충분히 평가할 수 없다. grasp attempt count, slip detection, lift success, placement success, recovery success를 분리해야 한다.

\section{Pushing, reaching, insertion}

Visuomotor policy는 grasping뿐 아니라 pushing, reaching, insertion에서도 사용된다. 각 task는 서로 다른 observation/action coupling을 가진다.

\begin{table}[ht]
\centering
\caption{Visuomotor task별 policy 요구사항}
\label{tab:visuomotor-tasks}
\small
\begin{tabularx}{\textwidth}{p{26mm}X X}
\toprule
Task & 필요한 관측 & 대표 action/failure \\
\midrule
Reaching & target pose, end-effector pose, obstacle & overshoot, depth error, calibration error \\
Grasping & object geometry, gripper state, contact cue & bad grasp point, slip, collision \\
Pushing & object pose, support friction, contact point & object rotation, loss of contact \\
Insertion & hole pose, peg pose, contact force, alignment & jamming, rim contact, excessive force \\
Wiping & surface normal, coverage state, pressure & missed area, too high/low pressure \\
\bottomrule
\end{tabularx}
\end{table}

이 표는 VLA에서도 그대로 중요하다. language instruction이 task를 지정할 수는 있지만, 각 task에서 필요한 visual feedback과 controller interface는 다르다.

\section{Action chunking and bimanual visuomotor policy}

고주파 control을 매 step 큰 모델로 생성하기 어렵기 때문에, policy가 짧은 horizon의 action chunk를 출력하는 구조가 많이 쓰인다.
\begin{equation}
  \bm{A}_{t:t+H}
  =
  \pi_{\theta}(\obs_{\le t},\conf_t,\task).
  \label{eq:visuomotor-action-chunk}
\end{equation}
실행 중에는 chunk의 앞부분만 쓰고 다음 관측에서 다시 chunk를 생성할 수 있다.
\begin{equation}
  \ctrlcmd_t
  =
  \controller(\bm{A}_{t:t+H}[0],\hat{\state}_t).
  \label{eq:chunk-execution}
\end{equation}

Bimanual manipulation에서는 action dimension이 커진다.
\begin{equation}
  \act_t
  =
  (\act^{left}_t,\act^{right}_t,\act^{sync}_t).
  \label{eq:bimanual-action}
\end{equation}
여기서 $\act^{sync}$는 두 arm 사이의 timing, relative pose, force balance를 나타낼 수 있다. cloth folding, cable manipulation, handover 같은 task에서는 이 synchronization이 핵심이다.

\section{From visuomotor policy to VLA}

VLA는 visuomotor policy에 language와 foundation-model representation을 추가한다.
\begin{equation}
  \act_t
  =
  \pi_{\theta}(\bm{I}_t,\conf_t,\task,\bel_t).
  \label{eq:vla-as-visuomotor}
\end{equation}
하지만 이 변화가 image-to-action 문제를 제거하지는 않는다. Language는 무엇을 해야 하는지 알려주지만, 어떻게 grasp할지, 어떤 frame에서 delta pose를 낼지, contact 후 어떤 correction을 할지, controller가 어떤 rate로 추적할지는 여전히 system design 문제이다.

\begin{table}[ht]
\centering
\caption{Visuomotor policy와 VLA의 차이}
\label{tab:visuomotor-vla}
\small
\begin{tabularx}{\textwidth}{p{30mm}X X}
\toprule
항목 & Visuomotor policy & VLA policy \\
\midrule
Task condition & task id, goal image, goal pose & language instruction, VLM embedding, subgoal \\
Observation & image, depth, proprioception & multimodal observation plus language/history \\
Action & delta pose, joint command, waypoint, chunk & token, continuous action, chunk, trajectory, skill call \\
Generalization & visual/task distribution 안의 일반화 & open-vocabulary semantics와 embodiment transfer 기대 \\
Remaining risk & distribution shift, calibration, contact, latency & 동일한 physical risk plus semantic ambiguity \\
\bottomrule
\end{tabularx}
\end{table}

\section{Evaluation metadata}

Visuomotor policy 또는 VLA policy를 평가할 때 다음 metadata가 필요하다.

\begin{itemize}
  \item camera setup: viewpoint, calibration, frame rate, latency
  \item proprioception and gripper state availability
  \item history length and memory architecture
  \item action type: delta pose, joint target, waypoint, chunk, token
  \item controller type and control frequency
  \item task-specific failure mode: grasp slip, insertion jamming, pushing loss of contact
  \item recovery policy and number of attempts
  \item train/test split: object, background, lighting, viewpoint, layout shift
\end{itemize}

\begin{casebox}{Case study: ACT/Diffusion Policy에서 VLA action chunk로}
ACT와 Diffusion Policy는 language 중심 VLA가 아니더라도 action chunk, multimodal action, temporal consistency가 robot policy에서 왜 중요한지 보여준다 \cite{act2023,diffusionpolicy2023}. VLA가 chunk나 diffusion/flow action을 쓰는 이유는 단순히 최신 생성모델을 적용하기 위해서가 아니라, visuomotor policy가 오래 겪어 온 multimodality와 timing 문제를 다루기 위해서다.
\end{casebox}

\chapterwork
{End-to-end visuomotor policy, visual foresight, Transporter, CLIPort, ACT/Diffusion Policy를 VLA의 실행 계보로 읽어라 \cite{levine2016gps,ebert2018visualforesight,florence2021transporter,cliport2022,act2023,diffusionpolicy2023}.}
{language가 추가되어도 image-to-action failure mode가 사라지지 않는 이유는 무엇인가?}
{하나의 pick-place task에 대해 camera setup, history length, action chunk horizon, controller rate, recovery label을 포함한 visuomotor evaluation sheet를 작성하라.}

\section{요약}

Visuomotor policy는 VLA의 이전 단계가 아니라 VLA의 실행 중심부에 남아 있는 문제다. Language와 VLM backbone은 task semantics와 open-vocabulary grounding을 제공하지만, robot action은 여전히 image, proprioception, history, controller interface, contact feedback 위에서 생성된다. 따라서 VLA를 이해하려면 transformer와 language model만 볼 수 없고, image-to-action policy가 real robot에서 어떤 실패를 겪어 왔는지 함께 이해해야 한다.

다음 장에서는 language와 vision representation 자체로 넘어간다. Transformer, VLM, grounding은 VLA가 open-vocabulary task를 이해하게 만든 기반이지만, 그 grounding이 robot action grounding과 어떻게 다른지 구분해야 한다.
